\documentclass[12pt]{article}
\usepackage[T1]{fontenc}
\usepackage[utf8]{inputenc}
\usepackage{lmodern}
\usepackage{textcomp}
\usepackage{enumitem}
\usepackage{geometry}
\usepackage{graphicx}
\usepackage{amsmath, amssymb}
\geometry{margin=1in}
\begin{document}
\begin{center}
Chemistry - Undergraduate Level Assessment 4 \
Total Marks: 40 \
Answer all questions.
\end{center}

\section*{Multiple Choice (10 marks)}
\begin{enumerate}[label=\arabic*.]
\item The anhydride of Ba(OH)2 is
\begin{enumerate}[label=(\alph*)]
    \item BaH 2
    \item BaOH
    \item Ba
    \item BaO
\end{enumerate}
\item Which of the following experimental observations were explained by Planck's quantum theory?
\begin{enumerate}[label=(\alph*)]
    \item Blackbody radiation curves
    \item Emission spectra of diatomic molecules
    \item Electron diffraction patterns
    \item Temperature dependence of reaction rates
\end{enumerate}
\item The solid-state structures of the principal allotropes of elemental boron are made up of which of the following structural units?
\begin{enumerate}[label=(\alph*)]
    \item B$_{12}$ icosahedra
    \item B$_{8}$ cubes
    \item B$_{6}$ octahedra
    \item B$_{4}$ tetrahedra
\end{enumerate}
\item Of the following atoms, which has the lowest electron affinity?
\begin{enumerate}[label=(\alph*)]
    \item F
    \item Si
    \item O
    \item Ca
\end{enumerate}
\item A set of hybrid $sp^3$ orbitals for a carbon atom is given above. Which of the following is NOT true about the orbitals?
\begin{enumerate}[label=(\alph*)]
    \item The orbitals are degenerate.
    \item The set of orbitals has a tetrahedral geometry.
    \item These orbitals are constructed from a linear combination of atomic orbitals.
    \item Each hybrid orbital may hold four electrons.
\end{enumerate}
\end{enumerate}

\section*{True / False (10 marks)}
\textit{Answer True or False}
\begin{enumerate}[label=\arabic*.]
\setcounter{enumi}{5}
\item True or False: In a three-component mixture, a maximum of five phases can coexist in equilibrium with one another.
\item True or False: The infrared absorption frequency of deuterium chloride (DCl) is shifted from that of hydrogen chloride (HCl) due to differences in reduced mass.
\item True or False: The NMR frequency of 31 P in a 20.0 T magnetic field is 345.0 MHz.
\item True or False: The orbital angular momentum quantum number, l, of the electron most easily removed from ground-state aluminum is 2.
\item True or False: Strong oxidizing conditions are necessary for a substance to exhibit paramagnetism and ferromagnetism.
\end{enumerate}

\section*{Long Form Response (20 marks)}
\textit{Answer each question with a clear and complete explanation.}
\begin{enumerate}[label=\arabic*.]
\setcounter{enumi}{10}
\item Discuss the characteristics that determine whether a compound behaves as a Lewis acid, and analyze why SCl 2 is the least likely to exhibit this behavior compared to other compounds.
\item Discuss the role of doping in semiconductors, specifically analyzing how arsenic-doped silicon functions as an n-type semiconductor and its implications for electronic applications.
\end{enumerate}

\vfill
\noindent\textit{End of Paper}
\end{document}